\chapter{Vorbereitung}

    \section{Kräfte zwischen Atomen}
        \subsection{Van-der-Waals Kräfte}

\begin{figure}[htb]
    \centering
    \includegraphics[width=0.6\textwidth]{Abb/vdw.pdf}
    \caption{Van-der-Waals Kraft zwischen zwei Teilchen}
    \label{vdw}
\end{figure}

Die Ladungsverteilung in Atomen ist nicht konstant, sondern unterliegt ständiger 
Fluktuation. Der Schwerpunkt der negativen Ladungen kann hierbei vom dem der
positiven Ladungen abweichen. Ist dies der Fall, so entsteht ein Dipol. 
Befindet sich nun ein zweites Teilchen in der Nähe dieses Atoms, so wird auch
in diesem ein Dipol induziert. Zeigt die positive Seite des ersten Atoms zu Atom 2,
so werden die Elektronen des zweiten Atoms angezogen. Ist es die negative Seite, so
werden die Elektronen abgestoßen. \par
Als Folge dessen synchronisieren sich die Ladungsänderungen der beiden Atome. Eine
schwache positive Anziehung ist die Folge. Diese ist proportional zu $\displaystyle
\frac{-1}{r^6}$.

        \subsection{Pauli-Abstoßung}

Nähern sich die Atome weiter an, so kommt es zu einem Überlappen der 
Elektronenorbitale. Das Pauli-Verbot verhindert hierbei, dass zwei Elektronen den
gleichen Zustand besetzen. Einige Elektronen werden folglich in einen energetisch
höheren Zustand gezwungen. \par
So führt eine Orbitalüberlagerung zu einer repulsiven Wechselwirkung. Die Kraft
ist proportional zu $\displaystyle \frac{1}{r^{12}}$.

        \subsection{Lennard-Jones Potential}

Bei sehr kleinen Abständen dominiert die Pauli-Abstoßung, bei größeren die 
van-der-Waals Kräfte. Die Summe aus beiden Potentialen wird Lennard-Jones 
Potential gennant. 
\[
   \phi (r) \propto \frac{A}{r^6} - \frac{B}{r^{12}}    
\]

\begin{figure}[hb]
    \centering
    \includegraphics[width=0.6\textwidth]{Abb/ljp.jpg}
    \caption{Das Lennard-Jones Potential als Summe der vdW-Wechselwirkung und
             der Pauli-Abstoßung}
    \label{ljp}
\end{figure}



Neben diesen Kräften können im Allgemeinen auch noch chemische Bindungskräfte, Kontaktkräfte, magnetische und elektrische Wechselwirkungen eine Rolle besitzen.
Bei unserem Aufbau spielen sie jedoch keine Rolle.



 \section{Das Atomic-Force-Mikroskop}
 
\begin{figure}[h]
    \centering
    \includegraphics[width=0.6\textwidth]{Abb/afm.jpg}
    \caption{Aufbau eines AFM}
    \label{afm}
\end{figure}
 
Mit dem Rasterkraftmikroskop, lassen sich nun die Abstoßung- bzw. Anziehungskräfte messen. Dazu fährt man mit dem allerwichtigsten Bestandteil, dem Cantilever, über die Probe. 
Dazu versetzt man den Cantilever in seine Resonanzfrequenz und beobachtet, wie stark diese Schwunung durch die Probe eingeschränkt wird.
Diese Amplitudenänderung lässt sich dann mittels Photodiode zu einem Kräfteprofil auslesen und in die Topographie der Probe umrechnen.
Der Cantilever ist ein schwingungsfähiger Balken, der eine pyramidale Spitze mit einer Dicke von nur wenigen Nanometern besitzt.
Er wird meist aus $Si_3N_4$ hergestellt und an der Spitze, durch Ätzung, zu einer sehr sehr dünnen Spitze geformt.
Seine Resonanzfrequenz befindet sich damit im kHz bis MHz Bereich.
Ein anderer wichtiger Bestandteil ist der Rastermechanismus, der die Spitze in die Nähe der Probe befördert. Die Spitze wird mittels Schrittmotoren auf einen ungefähren Abstand von einem Mikrometer angenähert und anschließend über piezoelektrische Stoffe weiter ausgerichtet.
Diese verändern ihre Größe, je nachdem wie viel Spannung an ihnen angelegt wird. 
Diese Vorrichtung führt letztendlich die Spitze über die Probe, in Dimensionen von 10 - 100 Mikrometern im Quadrat und einer Höhe von unter 10 Miktrometern.
Nun muss nur noch die Amplitudenänderung gemessen werden, dabei wird die Schingungsfrequenz konstant gehalten. 
Dann schießt man mit einem Laser auf eine polierte Rückseite des Cantilevers, welcher nun den Strahl auf eine Photodiode umlenkt. 
Damit können minimale Auslenkungen direkt druch elektrische Spannungsunterschiede aufgezeichnet werden.
Die Anziehenden und Abstoßenden Kräfte deformieren den Cantilever so, dass der Weg des Lasers im Angström Bereich umgelenkt wird. 
Eine Amplitudenänderungen von 70\% , zur Resonanzamplitude, ist maximal. 
Neben der Amplitudenmessung, ist es auch möglich die Frequenzänderung zu messen und die Möglichkeit des statischen Messens, also dem direkten Kontakt von Spitze und Probe.
Beide Methoden waren für die historische Entwicklung der Rasterkraftmikroskopie entscheidend, finden aber heutzutage keine Anwendung mehr.



